% GENERATED FILE. Edit canonical lessons, then run npm run generate:books.
% canonical-source-hash: fnv1a64:d16d4776a6f15ed7
% canonical-lessons: FR-C253-permis-de-conduire, FR-C253-casque, FR-C253-tunnel, FR-C253-chantier, FR-C253-immeuble

\chapter{Driving licence, Helmet, Tunnel, Building site, Block of flats}
\label{ch:fr-driving-licence-helmet-tunnel-building-site-block-of-flats}
\hlchaptermodality{253}

\begin{chapteropening}
\textbf{By the end of this chapter:} \emph{I can say a driving licence, a helmet, a tunnel, a building site and a block of flats.}

Complete the last lesson of chapter 253: I can say a driving licence, a helmet, a tunnel, a building site and a block of flats.
\end{chapteropening}

\section[le permis de conduire]{le permis de conduire — a driving licence}

\label{lesson:FR-C253-permis-de-conduire}



\begin{quote}
Before the new one: say the French for a stamp, then the French for a parcel.
\end{quote}

\subsection*{You'll want to know: le permis de conduire}
\textbf{le permis de conduire} — \textquotedblleft{}a driving licence\textquotedblright{}.

\begin{cousinweb}
\textbf{permis de conduire} — \emph{permis}, \textquotedblleft{}permitted\textquotedblright{}, from \emph{permettre}, \emph{de}, \textquotedblleft{}of\textquotedblright{}, and \emph{conduire}, \textquotedblleft{}to drive\textquotedblright{}: leave to drive.
\end{cousinweb}

\begin{grammarlens}[title={What you've built}]
One more word for this chapter.
\end{grammarlens}

\subsection*{Your turn}
\begin{itemize}
\raggedright
  \item \emph{Say it:} \emph{le permis de conduire}
  \item \emph{Say it:} \emph{le permis de conduire}, once more
  \item \emph{Say it:} say \emph{le permis de conduire}
  \item \emph{Recall:} say the French for a stamp, then the French for a parcel, then say \emph{le permis de conduire} again
\end{itemize}

\begin{tcolorbox}[breakable,colback=teal!4,colframe=teal!35!black,title={Before you move on}]
What does le permis de conduire mean? (\textquotedblleft{}A driving licence\textquotedblright{}.) Say it once more.
\end{tcolorbox}
\section[le casque]{le casque — a helmet}

\label{lesson:FR-C253-casque}



\begin{quote}
Before the new one: say the French for a parcel, then the French for a driving licence.
\end{quote}

\subsection*{You'll want to know: le casque}
\textbf{le casque} — \textquotedblleft{}a helmet\textquotedblright{}.

\begin{cousinweb}
\textbf{casque} $\leftarrow$ Spanish \emph{casco}, \textquotedblleft{}a skull, a helmet\textquotedblright{}, from \emph{cascar}, \textquotedblleft{}to break\textquotedblright{}, $\leftarrow$ Latin \emph{quassare}, \textquotedblleft{}to shatter\textquotedblright{} — borrowed. (EtymDB 2.1)
\end{cousinweb}

\begin{grammarlens}[title={What you've built}]
One more word for this chapter.
\end{grammarlens}

\subsection*{Your turn}
\begin{itemize}
\raggedright
  \item \emph{Say it:} \emph{le casque}
  \item \emph{Say it:} \emph{le casque}, once more
  \item \emph{Say it:} say \emph{le casque}
  \item \emph{Recall:} say the French for a parcel, then the French for a driving licence, then say \emph{le casque} again
\end{itemize}

\begin{tcolorbox}[breakable,colback=teal!4,colframe=teal!35!black,title={Before you move on}]
What does le casque mean? (\textquotedblleft{}A helmet\textquotedblright{}.) Say it once more.
\end{tcolorbox}
\section[le tunnel]{le tunnel — a tunnel}

\label{lesson:FR-C253-tunnel}



\begin{quote}
Before the new one: say the French for a driving licence, then the French for a helmet.
\end{quote}

\subsection*{You'll want to know: le tunnel}
\textbf{le tunnel} — \textquotedblleft{}a tunnel\textquotedblright{}.

\begin{cousinweb}
\textbf{tunnel} $\leftarrow$ English \emph{tunnel} (EtymDB 2.1) — borrowed back: English had taken it from Old French \emph{tonel}, \textquotedblleft{}a cask\textquotedblright{}.
\end{cousinweb}

\begin{grammarlens}[title={What you've built}]
One more word for this chapter.
\end{grammarlens}

\subsection*{Your turn}
\begin{itemize}
\raggedright
  \item \emph{Say it:} \emph{le tunnel}
  \item \emph{Say it:} \emph{le tunnel}, once more
  \item \emph{Say it:} say \emph{le tunnel}
  \item \emph{Recall:} say the French for a driving licence, then the French for a helmet, then say \emph{le tunnel} again
\end{itemize}

\begin{tcolorbox}[breakable,colback=teal!4,colframe=teal!35!black,title={Before you move on}]
What does le tunnel mean? (\textquotedblleft{}A tunnel\textquotedblright{}.) Say it once more.
\end{tcolorbox}
\section[le chantier]{le chantier — a building site}

\label{lesson:FR-C253-chantier}



\begin{quote}
Before the new one: say the French for a helmet, then the French for a tunnel.
\end{quote}

\subsection*{You'll want to know: le chantier}
\textbf{le chantier} — \textquotedblleft{}a building site\textquotedblright{}.

\begin{cousinweb}
\textbf{chantier} $\leftarrow$ Latin \emph{cantherius}, \textquotedblleft{}a gelding; a rafter, a trestle\textquotedblright{} (EtymDB 2.1) — inherited: first a timber support, then a work yard.
\end{cousinweb}

\begin{grammarlens}[title={What you've built}]
One more word for this chapter.
\end{grammarlens}

\subsection*{Your turn}
\begin{itemize}
\raggedright
  \item \emph{Say it:} \emph{le chantier}
  \item \emph{Say it:} \emph{le chantier}, once more
  \item \emph{Say it:} say \emph{le chantier}
  \item \emph{Recall:} say the French for a helmet, then the French for a tunnel, then say \emph{le chantier} again
\end{itemize}

\begin{tcolorbox}[breakable,colback=teal!4,colframe=teal!35!black,title={Before you move on}]
What does le chantier mean? (\textquotedblleft{}A building site\textquotedblright{}.) Say it once more.
\end{tcolorbox}
\section[l'immeuble]{l'immeuble — a block of flats}

\label{lesson:FR-C253-immeuble}



\begin{quote}
Before the new one: say the French for a tunnel, then the French for a building site.
\end{quote}

\subsection*{You'll want to know: l'immeuble}
\textbf{l'immeuble} — \textquotedblleft{}a block of flats\textquotedblright{}.

\begin{cousinweb}
\textbf{immeuble} $\leftarrow$ Latin \emph{immobilis}, \textquotedblleft{}that cannot be moved\textquotedblright{} (Lewis \& Short) — a learned borrowing: first property in land and buildings.
\end{cousinweb}

\begin{grammarlens}[title={What you've built}]
One more word for this chapter.
\end{grammarlens}

\subsection*{Your turn}
\begin{itemize}
\raggedright
  \item \emph{Say it:} \emph{l'immeuble}
  \item \emph{Say it:} \emph{l'immeuble}, once more
  \item \emph{Say it:} say \emph{l'immeuble}
  \item \emph{Recall:} say the French for a tunnel, then the French for a building site, then say \emph{l'immeuble} again
\end{itemize}

\begin{tcolorbox}[breakable,colback=teal!4,colframe=teal!35!black,title={Before you move on}]
What does l'immeuble mean? (\textquotedblleft{}A block of flats\textquotedblright{}.) Say it once more.
\end{tcolorbox}
